\input{../tex-inputs/latex-leseansicht-vorspann}

\section[Arthur Schnitzler an Gustav Schwarzkopf, {[}4?{]}. 7. 1900]{L04417 Arthur Schnitzler an Gustav Schwarzkopf, {[}4?{]}. 7. 1900}
\nopagebreak\mylabel{L04417v}
\rehead{ }\normalsize\beginnumbering\briefempfaengerindex{Schwarzkopf, Gustav@\textsc{Schwarzkopf, Gustav}!zzzSchnitzler, Arthur@\emph{von Arthur Schnitzler}!1900-07-042@{{[}4?{]}. 7. 1900}|(be}
\toendnotes[C]{\smallbreak\pagebreak[2]}
\correspDesc{Versand  durch Arthur Schnitzler am [4?]. 7. 1900 in Admont
\newline{}Übermittlung  am 5. 7. 1900 in Admont
\newline{}Zustellung  am 6. 7. 1900 in Wien 1/1 1
\newline{}Erhalt  durch Gustav Schwarzkopf am 6. 7. 1900 in Tiefer Graben 23}\toendnotes[C]{\smallbreak}
\Standort{DLA, A:Schnitzler, HS.1985.1.1897.}
\physDesc{Bildpostkarte, 172 Zeichen
\newline{}Handschrift: Bleistift, deutsche Kurrent
\newline{}Versand: 1) Stempel: »\nobreak{}\oindex{Admont@\textbf{Admont}|pwk}Admont, 5/7 {[}00{]}\nobreak{}«.   2) mit Bleistift von unbekannter Hand Ergänzung des Stockwerks (?):
                                    »IIr« bei der Adresse 3) Stempel: »\nobreak{}\oindex{I., Innere Stadt@\textbf{I., Innere Stadt}|pwk}Wien 1/1 1, 6. 7. \textcolor{gray}{00}, 8–\textcolor{gray}{9}½V., \textcolor{gray}{Bestellt}\nobreak{}«. }\toendnotes[C]{\smallbreak}\pstart{}{\pb}Hrn \textsc{Gustav
                  Schwarzkopf}\pend{}\pstart{}Wien\pend{}\pstart{}\textsc{I. Tiefer Graben
                     23}\oindex{Wien@\textbf{Wien}!I., Innere Stadt@\textbf{I., Innere Stadt}!Tiefer Graben 23@\textbf{Tiefer Graben 23}|pw}\pend{}{\bigskip}
\pstart
           \centering{}{\pb}\textcolor{gray}{\textbf{Gruss aus Admont\oindex{Admont@\textbf{Admont}|pw}.}}\pend
           \vspace{1em}
\pstart
           \noindent{}{\pb}Herzliche Grüße. Nove\uline{lle}\pwindex{Schnitzler, Arthur 15. 5. 1862 Wien – 21. 10. 1931 ebd.@\textsc{Schnitzler, Arthur} (15. 5. 1862 Wien – 21. 10. 1931 ebd.)!Frau Bertha Garlan. Roman@\strich\emph{Frau Bertha Garlan. Roman}|pwv} erhielten Sie wohl. \label{K_L04417-1v}\edtext{\textcolor{gray}{Freitag} dürft ich in Reichenau\oindex{Reichenau an der Rax@\textbf{Reichenau an der Rax}|pw}}{\lemma{\textnormal{\emph{Freitag … Reichenau}}}\Cendnote{\textnormal{Schnitzler reiste bereits am Donnerstag, dem 5. 7. 1900 nach Reichenau\oindex{Reichenau an der Rax@\textbf{Reichenau an der Rax}|pwk}. Da dies die Verwendung von ›heute‹
                  respektive ›morgen‹ in der Ausdrucksweise nahe legen würde, ist anzunehmen, dass
                  die Karte am Tag vor dem Poststempel und damit am Mittwoch abgefasst wurde.
               }}}\label{K_L04417-1}{ }ſein, dann hören Sie mehr von mir.\pend
           
\pstart
           Herzlichſt Ihr{\\[\baselineskip]}\spacefill\mbox{Arthur Sch}\pend
           \leftskip=0em{}\selectlanguage{ngerman}\endnumbering\briefempfaengerindex{Schwarzkopf, Gustav@\textsc{Schwarzkopf, Gustav}!zzzSchnitzler, Arthur@\emph{von Arthur Schnitzler}!1900-07-042@{{[}4?{]}. 7. 1900}|)be}\mylabel{L04417h}
\begin{anhang}
\end{anhang}\newcommand{\dateiname}{L04417}\newcommand{\titel}{Arthur Schnitzler an Gustav Schwarzkopf, [4?]. 7. 1900}\newcommand{\editorInnen}{Herausgegeben von Jahnke, SelmaMüller, Martin Anton}\input{../tex-inputs/latex-leseansicht-abspann}
